\clearpage
\section{Subgraphs and Order Embeddings (35 points)}

For this question, graph \(A\) is a subgraph of graph \(B\) when some induced
subgraph of \(B\) is graph-isomorphic to \(A\). Anchor nodes are not considered,
and every coordinate of an order embedding is non-negative.

The order-embedding constraint is
\[
  A\text{ is a subgraph of }B
  \quad\Longleftrightarrow\quad
  \vect{z}_A[i]\leq \vect{z}_B[i]
  \quad\text{for every dimension }i.
\]
Equivalently, write \(\vect{z}_A\preceq\vect{z}_B\) for coordinate-wise
inequality.

\subsection{Transitivity (8 points)}

Show that the subgraph relation is transitive: if graph \(A\) is a subgraph of
graph \(B\), and graph \(B\) is a subgraph of graph \(C\), then graph \(A\) is a
subgraph of graph \(C\).

Use the subgraph-isomorphism definition: if \(A\) is a subgraph of \(B\), then
there is a bijection \(f\) from \(V_A\) to a subset of \(V_B\) such that the
subgraph of \(B\) induced by
\[
  \{f(v):v\in V_A\}
\]
is graph-isomorphic to \(A\).

\begin{solution}
Write the three graphs as
\[
  A=(V_A,E_A),\qquad B=(V_B,E_B),\qquad C=(V_C,E_C).
\]
Because \(A\) is a subgraph of \(B\), there is a subset
\(U_B\subseteq V_B\) and a bijection
\[
  f:V_A\longrightarrow U_B
\]
such that \(f\) is a graph isomorphism from \(A\) to the induced subgraph
\(B[U_B]\).  In particular, for every \(u,v\in V_A\),
\[
  \{u,v\}\in E_A
  \quad\Longleftrightarrow\quad
  \{f(u),f(v)\}\in E_B.
  \tag{1}
\]
Similarly, because \(B\) is a subgraph of \(C\), there are a subset
\(U_C\subseteq V_C\) and a bijection
\[
  g:V_B\longrightarrow U_C
\]
such that \(g\) is a graph isomorphism from \(B\) to the induced subgraph
\(C[U_C]\).  Thus, for every \(x,y\in V_B\),
\[
  \{x,y\}\in E_B
  \quad\Longleftrightarrow\quad
  \{g(x),g(y)\}\in E_C.
  \tag{2}
\]
The right-hand side of (2) is exactly the edge relation in \(C[U_C]\),
because that subgraph is induced.
Now define
\[
  U_C'=g(U_B)\subseteq U_C
  \qquad\text{and}\qquad
  h=g\circ f:V_A\longrightarrow U_C'.
\]
Both \(f\) and \(g\) are bijections on their stated domains, so \(h\) is a
bijection from \(V_A\) onto \(U_C'\).  For any \(u,v\in V_A\), combining
(1) and (2) gives
\[
\begin{aligned}
  \{u,v\}\in E_A
  &\Longleftrightarrow \{f(u),f(v)\}\in E_B\\
  &\Longleftrightarrow \{g(f(u)),g(f(v))\}\in E_C\\
  &\Longleftrightarrow \{h(u),h(v)\}\in E\bigl(C[U_C']\bigr).
\end{aligned}
\]
The last equivalence holds because \(C[U_C']\) is the induced subgraph on
\(U_C'\): it contains precisely the edges of \(C\) whose two endpoints lie
in \(U_C'\).  Therefore \(h\) preserves both edges and non-edges and is a
graph isomorphism from \(A\) to \(C[U_C']\).
Since \(C[U_C']\) is an induced subgraph of \(C\), the definition of the
subgraph relation implies that \(A\) is a subgraph of \(C\).  Hence the
subgraph relation is transitive.
\end{solution}

\subsection{Anti-symmetry (8 points)}

Using the same subgraph-isomorphism definition, show that if graph \(A\) is a
subgraph of graph \(B\), and graph \(B\) is a subgraph of graph \(A\), then
\(A\) and \(B\) are graph-isomorphic.

\textbf{Hint.} What does the condition imply about the numbers of nodes in
\(A\) and \(B\), and why does that implication result in graph isomorphism?

\begin{solution}
Write
\[
  A=(V_A,E_A),\qquad B=(V_B,E_B),
\]
and assume, as in the homework setting, that the graphs are finite.  Since
\(A\) is a subgraph of \(B\), there are a subset
\(U_B\subseteq V_B\) and a bijection
\[
  f:V_A\longrightarrow U_B
\]
such that \(A\) is graph-isomorphic to the induced subgraph \(B[U_B]\).
Consequently,
\[
  |V_A|=|U_B|\leq |V_B|.
  \tag{1}
\]

Conversely, since \(B\) is a subgraph of \(A\), there are a subset
\(U_A\subseteq V_A\) and a bijection from \(V_B\) onto \(U_A\).  Hence
\[
  |V_B|=|U_A|\leq |V_A|.
  \tag{2}
\]
Combining (1) and (2) gives
\[
  |V_A|=|V_B|.
\]
It follows from \(U_B\subseteq V_B\) and
\(|U_B|=|V_A|=|V_B|\) that
\[
  U_B=V_B.
\]
Therefore the induced subgraph \(B[U_B]\) is the whole graph \(B\), not a
proper subgraph.

The original bijection \(f:V_A\to U_B\) is consequently a bijection
\(f:V_A\to V_B\).  Since it is an isomorphism from \(A\) to \(B[U_B]=B\),
for every \(u,v\in V_A\) it satisfies
\[
  \{u,v\}\in E_A
  \quad\Longleftrightarrow\quad
  \{f(u),f(v)\}\in E_B.
\]
Thus \(f\) preserves both edges and non-edges and is a graph isomorphism
between \(A\) and \(B\).  Hence the subgraph relation is anti-symmetric up to
graph isomorphism.
\end{solution}

\subsection{Common Subgraph (4 points)}

Consider a two-dimensional order-embedding space. Graphs \(A\) and \(B\) have
embeddings \(\vect{z}_A\) and \(\vect{z}_B\), and suppose the order-embedding
constraint is preserved perfectly. Prove that graph \(X\) is a common subgraph
of \(A\) and \(B\) if and only if
\[
  \vect{z}_X\preceq
  \min\{\vect{z}_A,\vect{z}_B\},
\]
where the minimum is taken coordinate-wise.

\begin{solution}
Let
\[
  \min\{\vect{z}_A,\vect{z}_B\}
  =
  \left(
    \min\{\vect{z}_A[0],\vect{z}_B[0]\},
    \min\{\vect{z}_A[1],\vect{z}_B[1]\}
  \right).
\]
By definition, \(X\) is a common subgraph of \(A\) and \(B\) exactly when
\(X\) is a subgraph of \(A\) and \(X\) is a subgraph of \(B\).

\textbf{(\(\Longrightarrow\))} Assume that \(X\) is a common subgraph of
\(A\) and \(B\).  Perfect preservation of the order-embedding constraint gives
\[
  \vect{z}_X\preceq\vect{z}_A
  \qquad\text{and}\qquad
  \vect{z}_X\preceq\vect{z}_B.
\]
Therefore, for each coordinate \(i\in\{0,1\}\),
\[
  \vect{z}_X[i]\leq \vect{z}_A[i]
  \quad\text{and}\quad
  \vect{z}_X[i]\leq \vect{z}_B[i].
\]
The number \(\vect{z}_X[i]\) is consequently no larger than the smaller of
\(\vect{z}_A[i]\) and \(\vect{z}_B[i]\):
\[
  \vect{z}_X[i]
  \leq
  \min\{\vect{z}_A[i],\vect{z}_B[i]\},
  \qquad i=0,1.
\]
These two coordinate inequalities are precisely
\[
  \vect{z}_X\preceq
  \min\{\vect{z}_A,\vect{z}_B\}.
\]

\textbf{(\(\Longleftarrow\))} Conversely, assume that
\[
  \vect{z}_X\preceq
  \min\{\vect{z}_A,\vect{z}_B\}.
\]
For each \(i\in\{0,1\}\), the definition of the minimum gives
\[
  \min\{\vect{z}_A[i],\vect{z}_B[i]\}
  \leq \vect{z}_A[i]
  \qquad\text{and}\qquad
  \min\{\vect{z}_A[i],\vect{z}_B[i]\}
  \leq \vect{z}_B[i].
\]
Combining these inequalities with the assumption yields
\[
  \vect{z}_X[i]\leq\vect{z}_A[i]
  \qquad\text{and}\qquad
  \vect{z}_X[i]\leq\vect{z}_B[i]
  \quad\text{for }i=0,1.
\]
Hence
\[
  \vect{z}_X\preceq\vect{z}_A
  \qquad\text{and}\qquad
  \vect{z}_X\preceq\vect{z}_B.
\]
Applying the reverse direction of the perfect order-embedding constraint,
\(X\) is a subgraph of both \(A\) and \(B\).  Thus \(X\) is a common
subgraph of \(A\) and \(B\).

Both implications hold, so
\[
  X\text{ is a common subgraph of }A\text{ and }B
  \quad\Longleftrightarrow\quad
  \vect{z}_X\preceq\min\{\vect{z}_A,\vect{z}_B\}.
\]
\end{solution}

\subsection{Incomparable Graphs (5 points)}

Suppose \(A\), \(B\), and \(C\) are pairwise non-isomorphic graphs, and none is
a subgraph of another. They are embedded into a two-dimensional order-embedding
space. Without loss of generality, assume
\[
  \vect{z}_A[0]>\vect{z}_B[0]>\vect{z}_C[0].
\]
What must this imply about the relationship among
\(\vect{z}_A[1]\), \(\vect{z}_B[1]\), and \(\vect{z}_C[1]\) if the
order-embedding constraint is satisfied perfectly?

\begin{solution}
The second coordinate must be ordered in the opposite direction:
\[
  \vect{z}_A[1]<\vect{z}_B[1]<\vect{z}_C[1].
\]
To see this, first compare \(A\) and \(B\).  We are given
\[
  \vect{z}_B[0]<\vect{z}_A[0],
\]
so the first coordinate already satisfies the inequality needed for
\(\vect{z}_B\preceq\vect{z}_A\).  If we also had
\(\vect{z}_B[1]\leq\vect{z}_A[1]\), then both coordinates would satisfy
\[
  \vect{z}_B[0]\leq\vect{z}_A[0]
  \qquad\text{and}\qquad
  \vect{z}_B[1]\leq\vect{z}_A[1].
\]
This would give \(\vect{z}_B\preceq\vect{z}_A\).  By the perfectly preserved
order-embedding constraint, \(B\) would then be a subgraph of \(A\), contrary
to the hypothesis.  Therefore
\[
  \vect{z}_B[1]>\vect{z}_A[1].
  \tag{1}
\]

Apply the same argument to \(B\) and \(C\).  Since
\(\vect{z}_C[0]<\vect{z}_B[0]\), the assumption
\(\vect{z}_C[1]\leq\vect{z}_B[1]\) would imply
\(\vect{z}_C\preceq\vect{z}_B\), and hence that \(C\) is a subgraph of \(B\).
This is also forbidden, so
\[
  \vect{z}_C[1]>\vect{z}_B[1].
  \tag{2}
\]
Combining (1) and (2) yields
\[
  \boxed{\vect{z}_A[1]<\vect{z}_B[1]<\vect{z}_C[1]}.
\]
The inequalities must be strict: equality in either comparison would still
make the vector with the smaller first coordinate coordinate-wise no larger
than the other vector, which would incorrectly assert a subgraph relation.
With this reverse ordering, every pair remains incomparable: the graph with
the larger first coordinate cannot be below the other one in dimension \(0\),
and the graph with the smaller first coordinate is above it in dimension
\(1\).
\end{solution}

\subsection{Limitation of a Two-Dimensional Order-Embedding Space (10 points)}

Suppose again that \(A\), \(B\), and \(C\) are pairwise non-isomorphic graphs,
and none is a subgraph of another. Construct graphs \(X\), \(Y\), and \(Z\),
each of which is a subgraph of one or more graphs in \(\{A,B,C\}\), so that
their embeddings must satisfy
\[
  \vect{z}_X\preceq\vect{z}_Y
  \qquad\text{and}\qquad
  \vect{z}_X\preceq\vect{z}_Z.
\]
You do not need to give embedding coordinates or draw the graphs. Specify only
which of \(A\), \(B\), and \(C\) contain \(X\), \(Y\), and \(Z\) as subgraphs,
and explain why the inequalities are guaranteed.

Use the following clarifications:
\begin{enumerate}
  \item Without loss of generality,
    \(\vect{z}_A[0]>\vect{z}_B[0]>\vect{z}_C[0]\).
  \item Each of \(X\), \(Y\), and \(Z\) may be a common subgraph of exactly one,
    two, or three graphs in \(\{A,B,C\}\).
  \item The construction must guarantee both inequalities whenever the
    order-embedding constraint is satisfied in two dimensions.
  \item It is not necessary to exhibit concrete graph drawings. The requested
    subgraph relationships and their justification are sufficient.
\end{enumerate}

\begin{solution}
We use the following containment pattern:
\begin{itemize}
  \item \(X\) is a common subgraph of all three graphs \(A,B,C\);
  \item \(Y\) is a subgraph of exactly \(B\), so \(Y\) is not a subgraph of
    \(A\) or \(C\);
  \item \(Z\) is also a subgraph of exactly \(B\), so \(Z\) is not a
    subgraph of \(A\) or \(C\).
\end{itemize}
The graphs \(Y\) and \(Z\) can be chosen to be different subgraphs of \(B\).
The important point is that their listed memberships are exact.

From the assumption in the question and the result of the preceding part,
the three parent embeddings have the strict ordering
\[
  \vect{z}_A[0]>\vect{z}_B[0]>\vect{z}_C[0]
  \qquad\text{and}\qquad
  \vect{z}_A[1]<\vect{z}_B[1]<\vect{z}_C[1].
  \tag{1}
\]

Because \(X\) is a subgraph of every parent, the order-embedding constraint
gives
\[
  \vect{z}_X\preceq\vect{z}_A,\qquad
  \vect{z}_X\preceq\vect{z}_B,\qquad
  \vect{z}_X\preceq\vect{z}_C.
\]
Taking the strongest bound in each coordinate and using (1),
\[
  \vect{z}_X[0]\leq\vect{z}_C[0]
  \qquad\text{and}\qquad
  \vect{z}_X[1]\leq\vect{z}_A[1].
  \tag{2}
\]

Consider \(Y\). Since \(Y\) is a subgraph of \(B\),
\[
  \vect{z}_Y[0]\leq\vect{z}_B[0],
  \qquad
  \vect{z}_Y[1]\leq\vect{z}_B[1].
  \tag{3}
\]
Moreover, \(Y\) is not a subgraph of \(A\). If
\(\vect{z}_Y[1]\leq\vect{z}_A[1]\), then (3) and (1) would also give
\(\vect{z}_Y[0]\leq\vect{z}_A[0]\), and hence
\(\vect{z}_Y\preceq\vect{z}_A\). Perfect preservation would then imply that
\(Y\) is a subgraph of \(A\), a contradiction. Therefore
\[
  \vect{z}_Y[1]>\vect{z}_A[1].
  \tag{4}
\]
Similarly, \(Y\) is not a subgraph of \(C\). The second inequality in (3)
and (1) already imply
\(\vect{z}_Y[1]\leq\vect{z}_B[1]<\vect{z}_C[1]\). Thus the only way that
\(\vect{z}_Y\not\preceq\vect{z}_C\) can hold is
\[
  \vect{z}_Y[0]>\vect{z}_C[0].
  \tag{5}
\]
Combining (2), (4), and (5) gives
\[
  \vect{z}_X[0]\leq\vect{z}_C[0]<\vect{z}_Y[0],
  \qquad
  \vect{z}_X[1]\leq\vect{z}_A[1]<\vect{z}_Y[1],
\]
so
\[
  \vect{z}_X\preceq\vect{z}_Y.
  \tag{6}
\]

Exactly the same argument applies to \(Z\): from
\(\vect{z}_Z\preceq\vect{z}_B\), together with
\(\vect{z}_Z\not\preceq\vect{z}_A\) and
\(\vect{z}_Z\not\preceq\vect{z}_C\), we
obtain
\[
  \vect{z}_Z[1]>\vect{z}_A[1]
  \qquad\text{and}\qquad
  \vect{z}_Z[0]>\vect{z}_C[0].
\]
Together with (2), this yields
\[
  \vect{z}_X[0]\leq\vect{z}_C[0]<\vect{z}_Z[0],
  \qquad
  \vect{z}_X[1]\leq\vect{z}_A[1]<\vect{z}_Z[1],
\]
and hence
\[
  \vect{z}_X\preceq\vect{z}_Z.
  \tag{7}
\]

Thus this containment pattern guarantees both required inequalities in every
perfect two-dimensional order embedding. The coordinate argument forces the
embedding to place \(X\) below both \(Y\) and \(Z\), even though the
construction only specifies their relationships with \(A,B,C\). We can
choose \(B\) to contain \(X,Y,Z\) as suitable induced subgraphs, while \(A\)
and \(C\) omit \(Y\) and \(Z\). We may additionally choose \(X\) not to be a
subgraph of either \(Y\) or \(Z\), illustrating the two-dimensional
representation limitation.
\end{solution}
